\subsection{محدودیت‌ها و تحلیل شکست}
تحلیل داده‌های به‌دست‌آمده ابعاد فنی مهمی از رفتار الگوریتم را نمایان می‌سازد:
\begin{itemize}
    \item \textbf{فلات دقت پایین (\lr{PASS\_LOW}) در مدل‌های پایه:} مقادیر دقت در مدل‌های \lr{DistilBERT} (\lr{50.92\%})، \lr{RoBERTa} (\lr{51.15\%}) و \lr{GPT-2} (\lr{50.92\%}) همگی در حوالی مرز پایه ۵۱ درصد قرار گرفته‌اند. علت اصلی این رخداد، تلفات اطلاعاتی ناشی از کسر میانگین در بردارهای کلیدی است؛ اگرچه صفرسازی میانگین سوگیری مقیاس را خنثی می‌کند، اما در غیاب یک تابع نگاشت بسامدی یا لایه‌های بهینه‌ساز پیشرفته، بخشی از انرژی کلیدها که حامل مفاهیم موقعیتی یا کلمات کلیدی پرتکرار متن هستند، از بین می‌رود و منجر به کاهش تمایز سیگنال به نویز در لایه‌های ابتدایی می‌گردد.
    \item \textbf{افزایش مصرف حافظه در مدل دیکودر علّی (\lr{GPT-2}):} ثبت حافظه \lr{1,221.30 MB} و کاهش توان به \lr{9,567.80 tok/s} در \lr{GPT-2} به دلیل ساختار تانسوری عملیات \lr{cumsum} بر روی ضرب خارجی تانسورهای سه‌بعدی ابعاد بالا است؛ این امر موجب تشکیل تانسورهای موقت بزرگ با اندازه $[B, H, N_k, d, d]$ در حین فراخوانی عملیات می‌شود که راندمان حافظه پردازنده گرافیکی را در طول‌های دنباله بلندتر تحت تاثیر قرار می‌دهد.
    \item \textbf{همگرایی در مدل \lr{Qwen}:} مدل \lr{Qwen} با دستیابی به دقت \lr{51.95\%} و وضعیت \lr{PASS}، توانایی سازگاری بدون خطای زمان اجرا با ساختار علّی و سرهای گروهی را به نمایش گذاشته است.
\end{itemize}
